\section{Irreducible relative cycle и вывод one-copy weight}

Предыдущий gate оставил архитектурный выбор между physical vectorlike
extension и classical mapping-cone sector. В classical branch единственной
невыведенной величиной была multiplicity \(k\) одинаковых relative copies,
дающая \(\lambda_{\rm rel}=k\). Настоящий gate проверяет, устраняет ли эту
свободу стандартный для finite geometry принцип irreducibility.

\subsection{Операторная система одной цепи}

Пусть
\begin{equation}
 \mathcal H_1=\bar4\oplus2_R\oplus4,
 \qquad
 \mathcal B_h=M_4(\mathbb C)\oplus M_2(\mathbb C)
 \,\oplus M_4(\mathbb C),
\end{equation}
а \(\mathcal D_\Delta\) --- self-adjoint chain operator с двумя ненулевыми
edges \(\Delta\) и \(\Delta^T\varepsilon_2\). Коммутант полной node algebra
состоит из block scalars:
\begin{equation}
 \mathcal B_h'
 =\{\alpha I_4\oplus\beta I_2\oplus\gamma I_4\}.
\end{equation}
Условие \([X,\mathcal D_\Delta]=0\) на первом edge даёт
\(\alpha=\beta\), а на втором --- \(\beta=\gamma\). Поэтому
\begin{equation}
 \boxed{\{\mathcal B_h,\mathcal D_\Delta\}'=\mathbb C I_{10}.}
 \label{eq:ps-relative-one-copy-commutant}
\end{equation}
Этот вывод требует только ненулевых edges и остаётся верным на rank-one
vacuum.

\subsection{Кратность одинаковых copies}

Для \(k\) coefficient-free одинаковых copies канонический lift имеет вид
\begin{equation}
 \mathcal H_k=\mathcal H_1\otimes\mathbb C^k,
 \qquad
 \mathcal B_{h,k}=\mathcal B_h\otimes I_k,
 \qquad
 \mathcal D_{\Delta,k}=\mathcal D_\Delta\otimes I_k.
\end{equation}
Тогда
\begin{equation}
 \boxed{
 \{\mathcal B_{h,k},\mathcal D_{\Delta,k}\}'
 =I_{10}\otimes M_k(\mathbb C),
 \qquad
 \dim_\mathbb C=k^2.}
 \label{eq:ps-relative-k-copy-commutant}
\end{equation}
Следовательно, lifted cycle irreducible ровно при \(k=1\).

\begin{theorem}[One-copy irreducibility selector]
В branch одинаковых relative copies без дополнительных copy-space operators
условие irreducibility однозначно выбирает
\begin{equation}
 k=1,
 \qquad
 \boxed{\lambda_{\rm rel}=1.}
\end{equation}
Тем самым общий weight determinant selector следует из connectedness,
common trace--Hodge metric и irreducibility, а не вводится как числовая
подгонка.
\end{theorem}

Machine classification для \(k=1,\ldots,5\) даёт commutant dimensions
\begin{equation}
 1,\ 4,\ 9,\ 16,\ 25,
\end{equation}
точно совпадающие с \(k^2\).

\subsection{Граница theorem: nontrivial copy dynamics}

При \(k>1\) можно разрушить \(M_k\)-commutant, если поместить на два edges
дополнительные noncommuting matrices в copy space. Например, для \(k=2\)
пара \(\sigma_x,\sigma_z\) имеет общий commutant \(\mathbb C I_2\). Но такая
конструкция уже не является multiplicity одинаковых copies: она добавляет
новый finite algebra/operator sector и новые orientation data.

\begin{nogocriterion}[Copy-mixing loophole]
Irreducibility не запрещает произвольную расширенную theory с noncommuting
copy-space dynamics. Она запрещает \(k>1\) именно в рассматриваемой
coefficient-free identical-copy branch. Любой \(k>1\) rescue должен быть
предъявлен как новый finite model и пройти отдельные order-one, reality,
Hessian и RG gates.
\end{nogocriterion}

\subsection{Полный \texorpdfstring{\(\Delta+C\)}{Delta+C} Hessian}

До устранения fixed-point auxiliary curvature действие равно
\begin{equation}
 V(\Delta,C)=-\rho^2+\tau^2
 +\|\mathcal D_\Delta^2-C\|_F^2,
 \qquad
 C\in\mathcal B_h^{\rm sa}.
 \label{eq:ps-relative-uneliminated-potential}
\end{equation}
Real dimension \(\mathcal B_h^{\rm sa}\) равна
\begin{equation}
 4^2+2^2+4^2=36.
\end{equation}
В stationary point
\(C_\star=E_h(\mathcal D_\Delta^2)\). Если \(J_E\) обозначает Jacobian
map \(\Delta\mapsto E_h(\mathcal D_\Delta^2)\), полный Hessian по
\(16+36=52\) real variables имеет block form
\begin{equation}
 H_{\rm full}=
 \begin{pmatrix}
 H_{\rm eff}+2J_E^TJ_E&-2J_E^T\\
 -2J_E&2I_{36}
 \end{pmatrix}.
 \label{eq:ps-relative-full-hessian-block}
\end{equation}
Его Schur complement по auxiliary block точно равен
\begin{equation}
 H_{\rm full}/(2I_{36})=H_{\rm eff}.
\end{equation}
Для \(\lambda_{\rm rel}=1\)
\begin{equation}
 \Spec H_{\rm eff}
 =\{8\sqrt2\ (1),\ 2\sqrt2\ (6),\ 0\ (9)\},
\end{equation}
а inertia полного Hessian равна
\begin{equation}
 \boxed{(43_+,9_0,0_-).}
 \label{eq:ps-relative-full-hessian-inertia}
\end{equation}
Numerical Schur-complement error равна
\(1.8\times10^{-15}\). Все девять zero modes уже присутствуют в effective
\(\Delta\)-sector; auxiliary block не создаёт новых flat или negative
directions.

\begin{theorem}[Classical branch conditional closure]
В irreducible coefficient-free mapping-cone branch multiplicity и общий
weight закрыты: \(k=\lambda_{\rm rel}=1\). Полный
\(\Delta+C\) Hessian не имеет negative modes. Остаётся встроить relative
cycle в полную KO6 real structure и вычислить mixed directions с
\(\phi\) и \(\Sigma_4\), которые не входят в
\eqref{eq:ps-relative-uneliminated-potential}.
\end{theorem}

\subsection{Литература и воспроизводимость}

Multiplicity matrices finite bimodules и их direct sums сверены с
T.~Krajewski, arXiv:hep-th/9701081. Использование irreducible minimal
Krajewski diagrams и потеря irreducibility при добавлении одинаковых family
copies сверены с J.-H.~Jureit and C.~A.~Stephan,
arXiv:hep-th/0610040. Irreducibility finite real triples и multiplicity-one
structure сопоставлены с A.~H.~Chamseddine and W.~D.~van Suijlekom,
arXiv:1904.12392. Relative spectral triple и doubling framework сверены с
I.~Forsyth, M.~Goffeng, B.~Mesland and A.~Rennie,
arXiv:1607.07143. Формулы
\eqref{eq:ps-relative-one-copy-commutant}--
\eqref{eq:ps-relative-full-hessian-inertia} являются project specialization.

Machine audit сохранён в
\path{s2t_v4_pati_salam_irreducible_relative_cycle_gate.py} и
\path{s2t_v4_pati_salam_irreducible_relative_cycle_gate_results.json}.